\chapter{TMDs and Global Fitting: Contact with Phenomenology}

Two interfaces connect the lattice program to the rest of particle
physics. Vertically, transverse-momentum-dependent distributions extend
the collinear picture and bring their own factorization physics ---
born in the paper below. Horizontally, global fits turn thousands of data
points into the PDF sets that lattice calculations benchmark against and
increasingly absorb as input. One landmark of each line is included here,
plus the TMD founding document.

% ----------------------------------------------------------------------
\section{Back-to-Back Jets: The Birth of TMD Physics}
\papercard{J.\ C.\ Collins, D.\ E.\ Soper}
{Back-to-back jets in QCD}
{Nucl.\ Phys.\ B \textbf{193}, 381 (1981)}
{DOI: 10.1016/0550-3213(81)90506-7 | companion: Collins--Soper--Sterman,
Nucl.\ Phys.\ B \textbf{250}, 199 (1985)}
\whyselected{The definitional paper of transverse-momentum-dependent
factorization ($\sim$4/99 explicit, universal within the library's TMD
line): every quasi-TMD article builds on its concepts.}
\physics{For color-singlet production with small transverse momentum,
collinear factorization fails; observables depend on a new class of
matrix elements with both light-cone and transverse separations. Collins
and Soper introduced the rapidity cut $\zeta$ to define these TMDs,
proved their factorization from hard, jet and soft sectors, and derived
the evolution equation in rapidity --- the Collins--Soper kernel ---
supplemented by DGLAP in $k_\perp$:
\begin{equation}
  \frac{\partial \ln F(x,\mathbf{k}_\perp;\zeta,\mu)}
       {\partial \ln\sqrt{\zeta}} = K(b_\perp;\mu),
  \qquad b_\perp \leftrightarrow \mathbf{k}_\perp .
\end{equation}
The companion CSS paper completed the proof for Drell--Yan. The kernel
$K$ is intrinsically nonperturbative at large $b_\perp$ --- which is
precisely where the lattice enters decades later.}
\impact{The library's quasi-TMD wave-function papers, soft-function
extractions and nonperturbative Collins--Soper-kernel determinations all
target this object; their renormalization puzzles (rapidity divergences
on the lattice) are direct descendants of the structure this paper
discovered.}
\cardend

% ----------------------------------------------------------------------
\section{CT18: The Hessian Standard}
\papercard{T.-J.\ Hou et al.\ (CTEQ-TEA Collaboration)}
{New CTEQ global analysis of quantum chromodynamics with high-precision
data from the LHC}
{Phys.\ Rev.\ D \textbf{103}, 014013 (2021)}
{arXiv:1912.10053 | DOI: 10.1103/PhysRevD.103.014013}
\whyselected{Original article of this library and its most-quoted
phenomenological benchmark ($\sim$6/99 among fit-corpus citations):
the default comparison set for nearly every lattice PDF extraction in the
corpus.}
\physics{CT18 evolves the CTEQ methodology to the full LHC era: Hessian
error analysis over an extended basis, NNLO theory with heavy-quark
scheme variants, and careful treatment of dataset tensions (notably the
ATLAS 7~TeV W/Z measurements, accommodated in the CT18A/CT18Z variants).
It delivers the standard-candle predictions (Higgs channels, W mass
sensitivity) against which both collider analyses and first-principles
inputs are judged. Its predecessor CT14 (Dulat et al., PRD \textbf{93},
033006 (2016), also an original article of this library) established the
parameterization framework; the Hessian method itself goes back to
Pumplin et al., JHEP \textbf{07}, 012 (2002).}
\impact{In this library CT18 appears as the reference curve in gluon-PDF,
quasi-PDF and DA comparisons; the ``lattice PDFs enter the global fit''
program is calibrated by asking how much CT-class uncertainties shrink
when lattice matrix elements join the data pool.}
\cardend

% ----------------------------------------------------------------------
\section{NNPDF3.1: The Monte Carlo School}
\papercard{R.\ D.\ Ball et al.\ (NNPDF Collaboration)}
{Parton distributions from high-precision collider data}
{Eur.\ Phys.\ J.\ C \textbf{77}, 663 (2017)}
{arXiv:1706.00428 | DOI: 10.1140/epjc/s10052-017-5199-5}
\whyselected{The Monte-Carlo-methodology counterpart of CT18
($\sim$8/99), original article of this library; its successor NNPDF4.0
(\emph{The Path to Proton Structure at One-Percent Accuracy}) is likewise
in the library and cites it as baseline.}
\physics{NNPDF represents PDFs as neural-network ensembles fitted by
Monte Carlo, with closure tests certifying that uncertainties neither
underestimate nor inflate truth. Version 3.1 added independent charm
(theory-motivated but fitted), top-differential and high-precision LHC
electroweak-boson data, reaching the accuracy regime where theory errors
(NNLO splitting functions, electroweak corrections) dominate. The
methodological chain --- closure tests, replica compression, positivity
and sum-rule enforcement --- became the field's quality standard and,
through NNPDF4.0's hyperparameter optimization, converged toward
machine-learning engineering itself.}
\impact{Together the two fitting schools bracket every lattice result of
this corpus: plots throughout the library overlay extractions on CT18
and NNPDF bands simultaneously. The deeper convergence --- global fitting
as machine learning, lattice sampling as generative modeling --- is the
subject of the next chapter.}
\cardend
